\chapter*{ЗАКЛЮЧЕНИЕ}
\addcontentsline{toc}{chapter}{ЗАКЛЮЧЕНИЕ}

Проведенный сравнительный анализ показал, что методы кластеризации существенно различаются по своим предположениям и области применимости. Алгоритмы разбиения на основе центров (k-средних, k-мод, k-прототипов) эффективны для больших наборов данных при заранее заданном числе кластеров и хорошо работают при компактной структуре кластеров. Метод k-медоидов отличается большей устойчивостью к выбросам и допускает использование произвольных мер расстояния. Нечеткий алгоритм c-средних и смешанная гауссовская модель целесообразны в задачах, где требуется учитывать степень принадлежности объектов к кластерам. Алгоритм DBSCAN подходит для обнаружения кластеров произвольной формы без задания их числа и выявления шума, тогда как иерархическая кластеризация удобна для анализа структуры данных и выбора уровня детализации разбиения. Таким образом, выбор метода кластеризации должен основываться на типе данных, предполагаемой форме кластеров и требованиях к интерпретации результата.

В ходе выполнения научно-исследовательской работы были решены следующие задачи:
\begin{itemize}
    \item[---] проведен анализ предметной области кластеризации данных;
    \item[---] рассмотрены существующие меры расстояний;
    \item[---] рассмотрены основные методы кластеризации;
    \item[---] выделены критерии сравнения методов кластеризации;
    \item[---] сравнены рассмотренные методы по выделенным критериям;
    \item[---] рассмотрены основные методы оценки качества кластеризации.
\end{itemize}

В рамках работы были выполнены все поставленные задачи. Цель работы была достигнута.